\section{Discussion and Limitations}
\label{sec:limitations}

The practical value of process control is selectivity with an inspectable
reason. In Study~A, \casepath removes false-positive changes while retaining
nearly as many correct changes as Direct. The benefit is largest on
procedural branches, where a familiar claimant document is a poor substitute
for the particular notice or confirmation that the rule requires. Study~B's
current-case scope intervention establishes a related design choice: an obligation
must inherit the conditions of its enclosing process, not merely satisfy its
own local condition. A system builder can inspect and test both dependencies
before asking a customer for another document.

The evidence has specific limits. Study~A tests known concepts in new
contexts and additions rather than justified withdrawals, uses unequal computation across pipelines, and does not meet its
preregistered broad-superiority gate. Its source-derived references contain
author judgements about which documents establish a fact, rather than expert
consensus; a deployed system would also need to select the governing policy
instead of combining rules from multiple insurers. Study~B tests complete
synthetic claims under matched budgets, but its native interface failure
leaves counterfactual graph correctness unvalidated. Its retrospective
current-case analysis isolates scope within dependent controls; the broader
request comparison measures realized output under failures.
Neither study measures customer effort, expert agreement, autonomous
acquisition or performance in another jurisdiction.

\section{Conclusion}
\label{sec:conclusion}

Evidence planning benefits from making the reason for a request executable.
\casepath represents the dependency from source condition to obligation,
fact, capability and document, and lets case state determine which demands
are active. The controlled comparison finds fewer unjustified checklist
changes; the current-case intervention shows that inherited scope removes
unjustified demand while retaining the same recovered obligations in development. These results support a concrete design
principle: derive evidence demands from explicit, case-conditioned
obligations, and validate both the requests and the process semantics that
produce them. The released benchmarks, controller, traces and offline
reproduction make that principle inspectable and testable.
